% ABOUTME: Related work section of the IEEE report (issue #8).
% ABOUTME: C-MAPSS, the capped RUL target, classical and deep FD001 models, and the gap tested here.
\section{Related Work}
\label{sec:related}

\subsection{Benchmark and target}
C-MAPSS simulates run-to-failure trajectories of aircraft turbofan engines
with modelled damage propagation~\cite{saxena2008damage}. Its FD001
subset is the simplest of the four: one operating condition, one fault mode
(high-pressure compressor degradation), and 100 training and 100 test
engines. Training engines run until failure; each test engine stops at an
unknown point before failure and its true RUL at that cycle is given
separately.

A linear RUL target implies that an engine degrades at a constant rate from
its first cycle, which the sensors do not show. Heimes~\cite{heimes2008rnn},
who estimated RUL with recurrent neural networks, proposed the
piecewise-linear target used in much of the later
work~\cite{babu2016cnn,ragab2020attention}: the target is held at a
constant maximum early in life and falls linearly to zero at failure. The
maximum is chosen from the data and differs between
subsets~\cite{babu2016cnn}; Ragab et al.\ use 125 cycles for FD001 and
FD003 and 130 for FD002 and FD004~\cite{ragab2020attention}. We cap
training targets at 125.

\subsection{Classical models}
Early data-driven methods for this benchmark include similarity-based
prognostics~\cite{wang2008similarity}. In the experiments of Babu et
al.~\cite{babu2016cnn}, support vector regression reached an RMSE of 20.96
cycles on FD001, ahead of relevance vector regression (23.80) and a
multilayer perceptron (37.56). Zhang et al.~\cite{zhang2017modbne} report
17.91 for a Random Forest and 15.67 for gradient boosting, against 15.04
for their ensemble of deep belief networks. These tree ensembles are the
closest published counterparts of our models.

\subsection{Deep sequence models}
Deep models take a window of normalised sensor readings as input and learn
their own features. Babu et al.~\cite{babu2016cnn} applied a convolutional
network along the time axis of 15-cycle windows (RMSE 18.45). Zheng et
al.~\cite{zheng2017lstm} used a long short-term memory
(LSTM)~\cite{hochreiter1997lstm} network, whose gated memory cells can carry
information over long time lags (16.14). Li et al.~\cite{li2018dcnn}
used a one-dimensional CNN with multiple input channels; at 12.61 it has
the lowest FD001 RMSE in the comparison of Ragab et
al.~\cite{ragab2020attention}. Ragab et al.\ add an attention mechanism
to an LSTM encoder-decoder over 30-cycle windows of 14 sensors and reach
12.63, with the lowest score $S$ (Section~\ref{sec:evaluation}) in that
comparison. Table~\ref{tab:related} summarises these methods;
Table~\ref{tab:published} in Section~\ref{sec:results} places our test
results among them.

\begin{table}[t]
\centering
\caption{Representative FD001 methods. RMSE values are those collected
in~\cite{ragab2020attention}, Table III, except SVR, which is
from~\cite{babu2016cnn}. Ours is scored against test labels clipped at 125,
as in Table~\ref{tab:published}. Cap is the maximum of the training
target. A dash marks a setting not stated in~\cite{babu2016cnn}
or~\cite{ragab2020attention}.}
\label{tab:related}
\footnotesize
\begin{tabular}{llrl}
\hline
Method & Input & RMSE & Cap \\
\hline
SVR \cite{babu2016cnn}                 & --                    & 20.96 & -- \\
Random Forest \cite{zhang2017modbne}   & --                    & 17.91 & -- \\
Grad.\ boosting \cite{zhang2017modbne} & --                    & 15.67 & -- \\
2D CNN \cite{babu2016cnn}              & sensors, 15 cycles    & 18.45 & yes \\
Deep LSTM \cite{zheng2017lstm}         & sensor sequence       & 16.14 & -- \\
1D CNN \cite{li2018dcnn}               & sensor windows        & 12.61 & -- \\
ATS2S \cite{ragab2020attention}        & 14 sensors, 30 cycles & 12.63 & 125 \\
XGBoost (ours)                         & engineered            & 14.08 & 125 \\
\hline
\end{tabular}
\end{table}

\subsection{Gap addressed}
The most accurate FD001 results come from models that learn features from
raw sensor windows. This work asks how far hand-built degradation features
take classical models. For each of the 14 informative sensors we compute
the rolling mean, standard deviation and least squares slope over 5, 10
and 20 cycles, after operating regime normalisation, and keep the
normalised readings, the cycle number and a regression-based health index
(Section~\ref{sec:features}). Elastic Net, Random Forest and XGBoost are
trained on these features and tuned with cross validation grouped by
engine, so that no engine appears in both the training and validation
folds and the validation error is not inflated by near-duplicate cycles
(Section~\ref{sec:evaluation}). Sequence models of the kind discussed
above are the subject of the second phase of the project.
